\documentclass{beamer}
\usefonttheme[onlymath]{serif}
\usepackage[T1]{fontenc}
\usepackage[utf8]{inputenc}
\usepackage[english]{babel}
\usepackage{amsmath}
\usepackage{amssymb}
\usepackage{amsthm}
\usepackage{gensymb}
\usepackage{parskip}
\usepackage{mathtools}
\usepackage{listings}
\usepackage{hyperref}
\usepackage{graphicx}
\usepackage{color}
\usepackage{enumerate}
\usepackage{verbatim}
\usepackage{minted}
\parskip 0pt


\DeclareMathOperator{\lcm}{lcm}
\newcommand\floor[1]{\left\lfloor#1\right\rfloor}
\newcommand\ceil[1]{\left\lceil#1\right\rceil}
\newcommand\abs[1]{\left|#1\right|}
\newcommand\p[1]{\left(#1\right)}
\newcommand\sqp[1]{\left[#1\right]}
\newcommand\cp[1]{\left\{#1\right\}}
\newcommand\norm[1]{\left\lVert#1\right\rVert}
\renewcommand\Im{\operatorname{Im}}
\renewcommand\Re{\operatorname{Re}}

\usetheme{metropolis}
\graphicspath{{../../shared/}}

\title{Complete search and greedy solutions}
\author{Atli FF}
\institute{\href{http://ru.is/td}{School of Computer Science} \\[2pt] \href{http://ru.is}{Reykjavík University}}
\titlegraphic{\hfill\includegraphics[height=0.6cm]{../shared/kattis}}
\date{\textbf{Árangursrík forritun og lausn verkefna}}

\begin{document}

\begin{frame}[plain]
    \titlepage
\end{frame}

\section*{Complete search / Brute force}

\begin{frame}[plain]
	\frametitle{Solution space}
	\begin{itemize}
		 \item The set of all possible solutions to a problem is called the \textit{solution space}
		 \item Note that this solution space will then contain all the wrong solutions too
		 \item Iterating over the entire solution space is called a complete search or brute force solution
	\end{itemize}
\end{frame}

\begin{frame}[plain]
	\frametitle{Iterating over solution spaces}
	\begin{itemize}
		 \item Iterating over different solution spaces is key to being able to brute force problems.
		 \item For the simplest cases, like ones you may have seen already in problems, we can just nest for/while loops.
	\end{itemize}
\end{frame}

\begin{frame}[plain]
	\frametitle{Iterating over a variable number of loops}
	\begin{itemize}
		 \item A common somewhat harder thing is iterating over a sequence of integers, with variable bounds
		 \item Say we want to test every vector of $n$ integers where each value can be from $0$ to $m - 1$.
		 \item We could do this usually with $n$ for loops, but if $n$ is an input this does not work
		 \item In python we can use itertools, C++ needs something handmade
	\end{itemize}
\end{frame}

\begin{frame}[plain, fragile]
	\frametitle{All vectors}
	\begin{minted}{cpp}
int n, m; cin >> n >> m;
vector<int> counter(n, 0);
while(true) {
  process_solution(counter); // whatever the problem needs
  bool done = true;
  for(int i = 0; i < n; ++i) {
    counter[i]++;
    if(counter[i] == m) counter[i] = 0;
    else {
      done = false;
      break;
    }
  }
  if(done) { break; }
}
	\end{minted}
\end{frame}

\begin{frame}[plain]
	\frametitle{Iterating over subsets}
	\begin{itemize}
		 \item For subsets we could do the same where we have a \texttt{vector<bool>} and index $i$ denotes whether we include the $i$-th element or not
		 \item But for this specific case we can do something much faster
		 \item \texttt{int}s consist of bits, so we can instead just have a number and let the $i$-th bit denote whether we include the $i$-th element or not
		 \item This means the subsets will just be the numbers from $0$ to $2^n - 1$ for $n$ elements, which can be iterated over with a simple for loop
	\end{itemize}
\end{frame}

\begin{frame}[plain]
	\frametitle{Bit masks}
	\begin{itemize}
		 \item These things are usually called bit masks
		 \item If $A = \{1, 2, 3, 4, 5, 6\}$ and $B = \{1, 3, 5, 6\}$ then the corresponding numbers in binary would be $111111$ and $110101$ since the first bit is the one we write last. In decimal these are 63 and 53.
		 \item This also gives us a whole load of useful operations that work on masks
	\end{itemize}
\end{frame}

\begin{frame}[plain]
	\frametitle{Bit operations}
	\begin{tabular}{|l|l|}
	\hline
	Code & Meaning \\
	\hline
	\texttt{0} & The empty set \\
	\texttt{(1 <{}< n) - 1} & The set of the first $n$ elements \\
	\texttt{1 <{}< k} & The set containing only the $k$-th element \\
	\texttt{A | B} & The union of $A$ and $B$ \\
	\texttt{A \& B} & The intersection of $A$ and $B$ \\
	\texttt{A \string^ B} & The symmetric difference of $A$ and $B$ \\
	\texttt{\string~ A} & The complement of $A$ \\
	\hline
	\end{tabular}
\end{frame}

\begin{frame}[plain, fragile]
	\frametitle{Iterating over permutations}
	\begin{itemize}
		 \item Say we want to iterate over all permutations of some vector/list.
		 \item Luckily this is built into a lot of languages:
		 \begin{itemize}
		 	\item \texttt{next\_permutation(v.begin(), v.end())} in C++
		 	\item \texttt{itertools.permutations} in Python
		 \end{itemize}
	\end{itemize}
	\begin{minted}{cpp}
int n = 5; vector<int> perm(n);
for(int i = 0; i < n; ++i) perm[i] = i + 1;
do {
	for(int i = 0; i < n; ++i) cout << perm[i] << ' ';
	cout << '\n';
} while(next_permutation(perm.begin(), perm.end()));
\end{minted}
\end{frame}

\begin{frame}[plain]
	\frametitle{Backtracking}
	\begin{itemize}
		 \item The methods to iterate over permutations and subsets were rather specialized
		 \item Backtracking is a general framework to iterate over complex spaces
		 \item Solves many classic problems like n-queens and sudoku
	\end{itemize}
\end{frame}

\begin{frame}[plain]
	\frametitle{Backtracking}
	\begin{itemize}
		 \item Define some initial "empty" state and have some notion of partial or complete states
		 \item For example in sudoku this is an empty grid, a partially filled grid and a fully numbered grid
		 \item Then define transitions to further states
		 \item In sudoku this would be filling in a number such that it doesn't create an immediate contradiction
	\end{itemize}
\end{frame}

\begin{frame}[plain]
	\frametitle{Backtracking}
	\begin{itemize}
		 \item Now start with your empty state
		 \item Use recursion to traverse all states by using the transitions
		 \item If the current state is invalid, stop exploring this branch
		 \item Process all complete states
	\end{itemize}
\end{frame}

\begin{frame}[plain, fragile]
	\frametitle{Backtracking - pseudo code}
	\begin{small}
	\begin{minted}{cpp}
state S;

void generate() {
	if(!is_valid(S)) return;
	
	if(is_complete(S)) print(S);
	
	for(each state P that S can transition to) {
		apply transition to P;
		generate();
		undo transition to P;
	}
}

S = empty state;
generate();
	\end{minted}
	\end{small}
\end{frame}

\begin{frame}[plain, fragile]
	\frametitle{Backtracking - Subsets}
	\begin{itemize}
		 \item We can even replicate earlier functionality this way
	\end{itemize}
	\begin{small}
	\begin{minted}{cpp}
const int n = 5; bool pick[n];

void generate(int index) {
  if(index == n) {
    for(int i = 0; i < n; ++i)
      if(pick[i]) cout << i << ' ';
    cout << '\n';
  } else {
    pick[index] = true;
    generate(index + 1); // pick element at index
    pick[index] = false;
    generate(index + 1); // don't pick element at index
  }
}
generate(0);
	\end{minted}
	\end{small}
\end{frame}

\begin{frame}[plain, fragile]
	\frametitle{Backtracking - Permutations}
	\begin{small}
	\begin{minted}{cpp}
const int n = 5; int perm[n]; bool used[n];

void generate(int index) {
  if(index == n) {
    for(int i = 0; i < n; ++i)
      cout << perm[i] + 1 << ' ';
    cout << '\n';
  } else {
    // decide what the element at index should be
    for(int i = 0; i < n; ++i) {
      if(!used[i]) {
        used[i] = true;
        perm[index] = i;
        generate(at + 1);
        used[i] = false; // remember to undo move!
      } } } }
memset(used, 0, n); generate(0);
	\end{minted}
	\end{small}
\end{frame}

\begin{frame}[plain]
	\frametitle{Backtracking - N queens}
	\begin{itemize}
		 \item Another classic backtracking problem is n-queens
		 \item We have a $n \times n$ chessboard and want to place $n$ queens on it so no two of them can attack each other
		 \item We could use bit tricks to iterate over all subsets of $n$ pieces in the board, but that would be too slow
		 \item Backtracking is much faster since we prune branches of computation early, it's almost universally good to do extra work earlier to prune branches when backtracking
	\end{itemize}
\end{frame}

\begin{frame}[plain, fragile]
	\frametitle{Backtracking - N queens}
	\begin{itemize}
		 \item We go through the cells in order
		 \item Our transition is placing a queen, or not placing a queen
		 \item We don't place a queen if it would be able to attack another placed queen
	\end{itemize}
	\begin{small}
	\begin{minted}{cpp}
const int n = 8;
bool has_queen[n][n];
int threatened[n][n];
int queens_left = n;

// generate function

memset(has_queen, 0, sizeof(has_queen));
memset(threatened, 0, sizeof(threatened));
generate(0, 0);
	\end{minted}
	\end{small}
\end{frame}

\begin{frame}[plain, fragile]
	\frametitle{Backtracking - N queens}
	\begin{scriptsize}
	\begin{minted}{cpp}
void generate(int x, int y) {
  if(y == n) generate(x + 1, 0); // move onto next column
  else if(x== n) { // we are at the end
    if(queens_left == 0) // this is a valid solution
      print(); // exact implementation not important
   } else {
     if(queens_left > 0 && threatened[x][y] == 0) {
       has_queen[x][y] = true;
       for(auto p : queen_threaten(x, y)) // good exercise to implement this!
         threatened[p.first][p.second]++; 
       queens_left--;
       generate(x, y + 1);
       has_queen[x][y] = false;  // now to undo the move
       for(auto p : queen_threaten(x, y))
         threatened[p.first][p.second]--; 
       queens_left++;
     }
     generate(x, y + 1); // also have to try leaving it empty
   }
}
	\end{minted}
	\end{scriptsize}
\end{frame}

\section*{Greedy algorithms}

\begin{frame}[plain]{Greedy algorithms}
    \begin{itemize}
        \item An algorithm that always makes \textit{locally} optimal moves is called greedy
        \item For some kinds of problems this will give a \textit{globally} optimal solution as well
        \item Seeing when this is the case can be very tricky, and if used in the wrong context the solution will get a WA verdict
    \end{itemize}
\end{frame}

\begin{frame}[plain]{Submitting greedy solutions}
    \begin{itemize}
        \item The tricky thing with these solutions are that it's often hard to know if you've made a mistake and thus get WA or if there's some hole in the greedy algorithm
        \item It's often easy to think of all kinds of greedy solutions, but they are very often wrong
        \item Generally one would like to consider complete search or dynamic programming first, but of course some problems do require greedy solutions
    \end{itemize}
\end{frame}

\begin{frame}[plain]{Coin change}
    \begin{itemize}
        \item A classical example is making change. Say you want to sum up $n$ and have only denominations of $1, 5$ and $10$, what's the least amount of coins you can give back?
        \item The greedy solution would be to just always give the biggest coin you can that's not too much. So for say $24$ we'd do $10, 10, 1, 1, 1, 1$.
        \item Is this always optimal?
    \end{itemize}
\end{frame}

\begin{frame}[plain]{Coin change}
    \begin{itemize}
        \item Well, it turns out to depend on the denominations. Say we have denominations of $1, 8$ and $20$.
        \item For $n = 24$ we then give back $20, 1, 1, 1, 1$ instead of the optimal $8, 8, 8$.
        \item We will come back to this problem tomorrow when we solve the general case using dynamic programming.
    \end{itemize}
\end{frame}

\begin{frame}[plain]{Taxi assignment}
    \begin{itemize}
        \item Let's consider another problem. You are managing a taxi company and today $n$ drivers showed up and you have $m$ cars. 
        \item But not all drivers and cars are created equal. Car $i$ has $h_i$ horsepower and driver $j$ can only handle at most $g_j$ horsepower.
        \item What's the greatest number of drivers you can pair to cars such that they can handle their car?
    \end{itemize}
\end{frame}

\begin{frame}[plain]{The greedy step}
    \begin{itemize}
        \item The greedy idea here is to simply pair each car to the worst driver that can still handle that car.
        \item Thus we start by sorting the drivers and cars and then simply linearly walk through each and pair them together.
        \item It might not be obvious, but this actually gives the best answer.
    \end{itemize}
\end{frame}

\begin{frame}[plain,fragile]{Implementation}
\begin{minted}{cpp}
int main() {
    int n, m; cin >> n >> m;
    vi a(n), b(m);
    for(int i = 0; i < n; ++i) cin >> a[i];
    for(int i = 0; i < m; ++i) cin >> b[i];
    sort(a.begin(), a.end());
    sort(b.begin(), b.end());
    int ans = 0;
    for(int i = 0, j = 0; i < m; ++i) {
        while(j < n && a[j] < b[i]) j++;
        if(j < n) ans++, j++;
    }
    cout << ans << '\n';
}
\end{minted}
\end{frame}

\begin{frame}[plain]{Sorting}
    \begin{itemize}
        \item Greedy algorithms very often involve sorting
        \item More generally they often involve always picking the ``extremal'' option out of the local options, in some sense
        \item Biggest, shortest, cheapest, first, etc.
    \end{itemize}
\end{frame}

\begin{frame}[plain]{Job scheduling}
    \begin{itemize}
        \item Say we have a list of jobs, each starting at some time $s_j$ and finishing at some time $f_j$
        \item What's the largest amount of jobs we can complete if they can't overlap?
    \end{itemize}
\end{frame}

\begin{frame}[plain]{Solution}
    \begin{itemize}
        \item The solution is shockingly simple, but not obviously correct
        \item Order the jobs by completion time $f_j$ and then walk through them
        \item If you can complete a job in addition to the ones you've already picked, pick it
        \item The jobs you've picked by the end are the solution
    \end{itemize}
\end{frame}

\begin{frame}[plain]{Proof of correctness}
    \begin{itemize}
        \item Why is this correct though? Let's prove it.
        \item Suppose the algorithm is not optimal. Say we pick jobs of indices $i_1, i_2, \dots, i_k$ but a better solution picks $j_1, j_2, \dots, j_l$.
        \item Say the solutions agree on the first $r$ jobs (possibly $0$). 
        \item Now neither $i_{r+1}$ nor $j_{r+1}$ clash with the jobs $i_1 = j_1, i_2 = j_2, \dots, i_r = j_r$. But because we ordered things by end time, we must have that job $i_{r+1}$ ends no later than $j_{r+1}$. But then we could just as well have picked $i_{r+1}$. But this holds for any $r$, so by induction we have that $i_1, \dots, i_k$ is no worse than $j_1, \dots, j_l$, which gives a contradiction.
        \item Thus the algorithm is optimal.
    \end{itemize}
\end{frame}

\end{document}

\end{document}

